\section{Introduction}
\label{sec:intro}

Statistical arbitrage covers trading strategies that exploit predictable, short-lived
deviations in the relative prices of securities while holding little exposure to the market
as a whole. The simplest form is pairs trading \citep{gatev2006}: find two stocks that have
historically moved together, and when the spread between them widens, buy the cheap one and
short the expensive one. Generalising from pairs to a whole cross-section requires a model of
what ``moving together'' means for hundreds of stocks at once. A factor model provides one.

\citet{avellaneda2010} (hereafter AL) proposed using principal component analysis (PCA) of the
empirical correlation matrix to extract the factors directly from the data. Each stock's return
is decomposed into a systematic part, explained by a small number of \emph{eigenportfolios},
and an idiosyncratic residual. The cumulative residual is modelled as a mean-reverting
Ornstein--Uhlenbeck (OU) process. When a stock's residual is far below its estimated
equilibrium, the stock is bought and its factor exposure hedged away. Symmetrically, stocks with
residuals far above equilibrium are shorted. The resulting portfolio is close to market- and
factor-neutral, so its profits come from the mean reversion of the residuals.

AL report attractive risk-adjusted returns for US equities over 1997--2007 and note a decline in
performance in the later years of their sample. This project asks three questions:
\begin{enumerate}
  \item Does the PCA residual signal still carry information in the S\&P~500 after 2010?
  \item If it does, is the strategy profitable after realistic trading costs, and if not, where
        do the costs come from?
  \item Which design choices reduce those costs without destroying the signal, and do the
        improvements survive a held-out test period?
\end{enumerate}

\paragraph{Contributions.}
\begin{itemize}
  \item An open, tested implementation of the AL strategy with a pluggable data layer, so the
        same code can run on free data (Yahoo Finance) or survivorship-free data (CRSP).
  \item A decomposition of the strategy's turnover. It shows that the cost of hedging, which is
        incurred mainly when positions open and close, is as large as the cost of trading the
        signal itself.
  \item Evidence that hedging fewer factors, chosen in line with random-matrix theory, and a
        no-trade band roughly double the edge per unit traded.
  \item Evidence that a large part of the measured edge depends on trading at the signal's own
        closing price: a one-day execution lag removes up to half of the gross Sharpe ratio.
  \item A deep-learning signal extractor, trained on the net-of-cost Sharpe ratio in a strict
        walk-forward, that is far more robust to execution lag than the rule-based strategy, and a
        decomposition showing how much of its edge is reversal and how much a systematic tilt.
  \item An evaluation protocol with a strict in-sample/test split enforced by the tooling, and an
        explicit count of every variant tried, so the final results can be corrected for
        multiple testing.
\end{itemize}

\paragraph{Outline.}
\Cref{sec:related} reviews related work. \Cref{sec:data} describes the data,
\cref{sec:method} the methodology, and \cref{sec:impl} the implementation and its validation.
\Cref{sec:protocol} sets out the evaluation protocol, \cref{sec:results} presents results, and
\cref{sec:discussion} discusses them and their limitations.
